\subsection{Bredhja e rastit dhe kufiri i difuzionit}
Në një përmasë, hapi $\xi_n$ merr vlerat $\pm\ell$ me probabilitete të barabarta dhe
\begin{equation}
 X_N=\sum_{n=1}^{N}\xi_n.
\end{equation}
Pavarësia jep $\E[X_N]=0$ dhe $\Var(X_N)=N\ell^2$. Nëse një hap zgjat $\Delta t$, atëherë $t=N\Delta t$ dhe
\begin{equation}
 \E[X(t)^2]=\frac{\ell^2}{\Delta t}t=2Dt,
 \qquad D=\frac{\ell^2}{2\Delta t}.
\end{equation}
Për të marrë kufirin e vazhdueshëm mbahen $D$ dhe $t$ të fiksuara ndërsa $\ell,\Delta t\to0$. Shpërndarja e pozicionit i afrohet
\begin{equation}
 p(x,t)=\frac{1}{\sqrt{4\pi Dt}}\exp\left(-\frac{x^2}{4Dt}\right),
\end{equation}
që zgjidh $p_t=Dp_{xx}$. Ky është një shembull themelor ku një model mikroskopik stokastik çon në ekuacion makroskopik determinist.

Me drift, probabilitetet nuk janë të barabarta. Në kufirin e duhur shfaqet ekuacioni adveksion--difuzion
\begin{equation}
 p_t=-v p_x+Dp_{xx}.
\end{equation}
Kushtet thithëse përfaqësojnë largimin e grimcës; kushtet reflektuese ruajnë probabilitetin. Koha e parë e kalimit është ndryshore rasti dhe trajektoret që nuk e arrijnë kufirin brenda kohës së simulimit janë të censuruara, jo zero.

\subsection{Zinxhirët Markov dhe shpërndarja stacionare}
Për një zinxhir me hapësirë të fundme dhe matricë kalimi $P$, shpërndarja rresht evoluon si
\begin{equation}
 \vect\pi_{n+1}^T=\vect\pi_n^TP.
\end{equation}
Shpërndarja stacionare plotëson $\vect\pi^TP=\vect\pi^T$, $\pi_i\ge0$ dhe $\sum_i\pi_i=1$. Nëse zinxhiri është i pakthyeshëm dhe aperiodik, shpërndarja konvergon te stacionarja pavarësisht gjendjes fillestare. Shpejtësia lidhet me boshllëkun spektral ndërmjet vlerës vetjake $1$ dhe vlerës së dytë në modul.

Balanca e detajuar,
\begin{equation}
 \pi_iP_{ij}=\pi_jP_{ji},
\end{equation}
është kusht i mjaftueshëm për stacionaritet, sepse mbledhja sipas $i$ jep $(\vect\pi^TP)_j=\pi_j$. Ai nuk është kusht i domosdoshëm: mund të ketë rryma stacionare që nuk anulohen çift më çift.

\subsection{Autokorrelacioni dhe madhësia efektive}
Për një seri stacionare $A_n$, autokovarianca është
\begin{equation}
 C(k)=\E[(A_n-\mu)(A_{n+k}-\mu)],
 \qquad \rho(k)=C(k)/C(0).
\end{equation}
Varianca e mesatares për $N$ vëzhgime është
\begin{align}
 \Var(\overline A)&=\frac{1}{N^2}\sum_{i,j}\Cov(A_i,A_j)\\
 &\simeq\frac{C(0)}{N}\left[1+2\sum_{k=1}^{\infty}\rho(k)\right]
 =\frac{2\tau_{\mathrm{int}}C(0)}{N},
\end{align}
ku
\begin{equation}
 \tau_{\mathrm{int}}=\frac12+\sum_{k=1}^{\infty}\rho(k).
\end{equation}
Prandaj
\begin{equation}
 N_{\mathrm{eff}}=\frac{N}{2\tau_{\mathrm{int}}}.
\end{equation}
Formula $s/\sqrt N$ nënvlerëson gabimin kur mostrat janë të korreluara. Shuma e autokorrelacionit duhet prerë në një vonesë ku zhurma e vlerësimit nuk dominon; alternativa praktike janë mesataret në blloqe dhe shumë zinxhirë të pavarur.

\subsection{Koha e përzierjes dhe diagnostika}
Koha e përzierjes mat sa shpejt shpërndarja e zinxhirit harron gjendjen fillestare. Heqja e një numri fiks hapash pa diagnostikë nuk garanton termalizim. Vëzhgohen trajektoret, mesataret kumulative dhe krahasimi i zinxhirëve nga gjendje të ndryshme. Hollimi i zinxhirit zakonisht hedh informacion; ruajtja e të gjitha mostrave dhe përdorimi i $N_{\mathrm{eff}}$ është më transparent, përveç kufizimeve të ruajtjes ose kostos së vlerësimit.

\subsection{Përmbledhje dhe ushtrime}
Bredhja e rastit lidh fluktuacionet mikroskopike me difuzionin. Zinxhiri Markov shton varësinë kohore dhe kërkon korrigjimin e gabimit standard me autokorrelacionin.
\begin{enumerate}
\item Nxirrni $\E[r^2]=2dDt$ për bredhjen në $d$ përmasa.
\item Zgjidhni kohën mesatare të daljes nga një interval me kufij thithës.
\item Gjeni shpërndarjen stacionare dhe vlerat vetjake të një matrice kalimi $3\times3$.
\item Simuloni një zinxhir me dy gjendje dhe krahasoni $N_{\mathrm{eff}}$ teorik me atë të vlerësuar.
\item Tregoni me shembull se stacionariteti nuk nënkupton balancë të detajuar.
\end{enumerate}
